% GENERATED FILE. Edit canonical lessons, then run npm run generate:books.
% canonical-source-hash: fnv1a64:990a97baedc1e40d
% canonical-lessons: HI-C92-thoda, HI-C92-achchha-buraa, HI-C92-se-comparison, HI-C92-kitna, HI-C92-kitna-achchha, HI-C92-repaso-maatra, HI-C92-sintesis-maatra

\chapter{How Much, How Well, How Big}
\label{ch:hi-maatra}
\hlchaptermodality{100}

\begin{chapteropening}
\textbf{By the end of this chapter:} \emph{I can hedge an answer, say that something was done well or badly, compare two things, ask a quantity and pay a compliment.}

Answer the listening item the first practice paper sets -- haan, thodii -- and say why the ending rather than the word is what the item tests.
\end{chapteropening}

\section[thoṛā]{\hi{थोड़ा} — the honest yes}

\label{lesson:HI-C92-thoda}



\begin{quote}
Somebody asks whether you speak Hindi. You can say yes and you can say no. Try to say \textbf{a little}, which is the true answer, and find you cannot.
\end{quote}

\subsection*{You'll want to know: \hi{थोड़ा}}
\textbf{\hi{थोड़ा}} (\emph{thoṛā}) — a little, a bit, some.

\textbf{\hi{हाँ}, \hi{थोड़ी।}} — Yes, a little. That two-word answer is a listening item on a practice paper.

\begin{cousinweb}
It ends in \textbf{-\hi{आ}}, so it bends like every other word that does:

\noindent
\begin{tabularx}{\linewidth}{@{}>{\raggedright\arraybackslash}X>{\raggedright\arraybackslash}X@{}}
\toprule
\textbf{what is being hedged} & \textbf{the form} \\
\midrule
a masculine thing & \hi{थोड़ा} \\
a feminine thing & \hi{थोड़ी} \\
\bottomrule
\end{tabularx}

That is why the paper's answer is \textbf{\hi{थोड़ी}} rather than \textbf{\hi{थोड़ा}} — the question was about \textbf{\hi{हिंदी}}, and \emph{hindī} is feminine. The hedge agrees with what it hedges, not with the person speaking.

This catches English speakers out, because \emph{a little} in English is a fixed lump that never changes. Here it is an adjective, and adjectives agree.
\end{cousinweb}

\begin{grammarlens}[title={Grammar Lens: where it stands}]
\noindent
\begin{tabularx}{\linewidth}{@{}>{\raggedright\arraybackslash}X>{\raggedright\arraybackslash}X@{}}
\toprule
\textbf{sentence} & \textbf{meaning} \\
\midrule
\hi{मैं} \hi{थोड़ी} \hi{हिंदी} \hi{बोलता} \hi{हूँ।} & I speak a little Hindi. \\
\hi{हाँ}, \hi{थोड़ी।} & Yes, a little. \\
\bottomrule
\end{tabularx}

In front of the noun, like any adjective. And on its own as a complete answer, where the noun it agrees with is the one the question named.

A hedge is worth more at this level than it looks. \textbf{It is the difference between claiming fluency and being believed}, and an examiner hears the difference.
\end{grammarlens}

\subsection*{Your turn}
Say these aloud:

\begin{itemize}
\raggedright
  \item \textquotedblleft{}\hi{थोड़ा}\textquotedblright{} and then \textquotedblleft{}\hi{थोड़ी}\textquotedblright{}
  \item \textquotedblleft{}\hi{मैं} \hi{थोड़ी} \hi{हिंदी} \hi{बोलता} \hi{हूँ}\textquotedblright{}
  \item \textquotedblleft{}\hi{हाँ}, \hi{थोड़ी}\textquotedblright{} — and say why it is not \textbf{\hi{थोड़ा}}
  \item the same answer as a woman would say the whole sentence
\end{itemize}

\begin{tcolorbox}[breakable,colback=teal!4,colframe=teal!35!black,title={Before you move on}]
Why is the paper's answer \textbf{\hi{थोड़ी}}? (\emph{Hindī} is feminine, and the hedge agrees with what it hedges.) Does English change \emph{a little} the same way? (No — it is a fixed lump there.)
\end{tcolorbox}
\section[acchā / burā]{\hi{अच्छा} / \hi{बुरा} — well and badly, from words you already own}

\label{lesson:HI-C92-achchha-buraa}



\begin{quote}
You have \textbf{\hi{अच्छा}} for \emph{good} and \textbf{\hi{बुरा}} for \emph{bad}. English needs two different words to say \emph{well} and \emph{badly}. Guess what Hindi needs.
\end{quote}

\subsection*{You'll want to know: the manner use}
\textbf{\hi{अच्छा}} is also \emph{well}. \textbf{\hi{बुरा}} is also \emph{badly}.

\textbf{\hi{वह} \hi{अच्छा} \hi{बोलता} \hi{है।}} — He speaks well. \textbf{\hi{वह} \hi{बुरा} \hi{बोलता} \hi{है।}} — He speaks badly.

\begin{cousinweb}
English keeps two shelves of words — \emph{good / well}, \emph{bad / badly} — and puts \textbf{-ly} on the second shelf. Hindi keeps one. The same word describes a thing or describes an action, and its position in the sentence says which:

\noindent
\begin{tabularx}{\linewidth}{@{}>{\raggedright\arraybackslash}X>{\raggedright\arraybackslash}X@{}}
\toprule
\textbf{where it stands} & \textbf{what it describes} \\
\midrule
beside a noun & the thing — a \textbf{good} book \\
in front of a verb & the action — speaks \textbf{well} \\
\bottomrule
\end{tabularx}

So \emph{he is good} and \emph{he speaks well} use the same word in Hindi, and an English speaker's instinct to reach for a second form is the thing to unlearn.

\textbf{\hi{अच्छा}} is also where you can now read the letter \textbf{\hi{छ}}, which arrived a moment ago in the alphabet. Find it in the middle of the word.
\end{cousinweb}

\begin{grammarlens}[title={Grammar Lens: it still agrees}]
\noindent
\begin{tabularx}{\linewidth}{@{}>{\raggedright\arraybackslash}X>{\raggedright\arraybackslash}X@{}}
\toprule
\textbf{speaker} & \textbf{sentence} \\
\midrule
a man & \hi{वह} \hi{अच्छा} \hi{बोलता} \hi{है।} \\
a woman & \hi{वह} \hi{अच्छी} \hi{बोलती} \hi{है।} \\
\bottomrule
\end{tabularx}

Even doing the work of an English adverb, it is still an adjective, so it still bends. English adverbs never do, which is the second half of the same instinct to unlearn.

The ordinary spoken way to say \emph{well} is also \textbf{\hi{अच्छी} \hi{तरह}}, \textquotedblleft{}in a good manner\textquotedblright{} — useful to recognise, and the bare adjective is enough to produce.
\end{grammarlens}

\subsection*{Your turn}
\begin{itemize}
\raggedright
  \item \emph{Say it:} \textquotedblleft{}\hi{वह} \hi{अच्छा} \hi{बोलता} \hi{है}\textquotedblright{} — he speaks well
  \item \emph{Say it:} \textquotedblleft{}\hi{वह} \hi{बुरा} \hi{बोलता} \hi{है}\textquotedblright{} — he speaks badly
  \item \emph{Say it:} the same pair as a woman would say them
  \item \emph{Look:} at \hi{अच्छा} and put your finger on the letter \textbf{\hi{छ}}
\end{itemize}

\begin{tcolorbox}[breakable,colback=teal!4,colframe=teal!35!black,title={Before you move on}]
How many words does Hindi need for \emph{good} and \emph{well}? (One.) What tells you which job it is doing? (Where it stands.) Does it still bend for gender when it means \emph{well}? (Yes — it is still an adjective.)
\end{tcolorbox}
\section[se baṛā]{\hi{से} \hi{बड़ा} — bigger, with no word for \emph{-er}}

\label{lesson:HI-C92-se-comparison}



\begin{quote}
You have \textbf{\hi{बड़ा}} and \textbf{\hi{छोटा}}. Try to say that one thing is \emph{bigger than} another, and notice English changed the adjective to do it.
\end{quote}

\subsection*{You'll want to know: \hi{से} \hi{बड़ा}}
\textbf{\hi{आदमी} \hi{बच्चे} \hi{से} \hi{बड़ा} \hi{है।}} — The man is bigger than the child.

The adjective \textbf{does not change}. The thing compared against takes \textbf{\hi{से}}.

\begin{cousinweb}
English has two machines for this — \textbf{-er} on short adjectives and \emph{more} in front of long ones — and both are extra grammar. Hindi has neither:

\noindent
\begin{tabularx}{\linewidth}{@{}>{\raggedright\arraybackslash}X>{\raggedright\arraybackslash}X@{}}
\toprule
\textbf{English} & \textbf{Hindi} \\
\midrule
bigger \textbf{than} X & X \textbf{\hi{से}} \hi{बड़ा} \\
smaller \textbf{than} X & X \textbf{\hi{से}} \hi{छोटा} \\
\bottomrule
\end{tabularx}

The postposition you learned for \emph{from} and \emph{by} is doing this too, and the logic holds: the comparison \textbf{proceeds from} the thing measured against. Measured \textbf{from} that book, this one is big.

So there is no comparative form to learn in this language. \textbf{There is one more use of a word you already have.}
\end{cousinweb}

\begin{grammarlens}[title={Grammar Lens: word order, and the reason the note was wrong}]
\noindent
\begin{tabularx}{\linewidth}{@{}>{\raggedright\arraybackslash}X>{\raggedright\arraybackslash}X@{}}
\toprule
\textbf{slot} & \textbf{what goes in it} \\
\midrule
first & the thing being described \\
second & what it is measured against, plus \textbf{\hi{से}} \\
third & the adjective, then the verb \\
\bottomrule
\end{tabularx}

Note what \textbf{\hi{बच्चे} \hi{से}} is doing there: that is the oblique from the last chapter, earning its keep the very next time a postposition turns up.

English puts \emph{than X} \textbf{after} the adjective; Hindi puts it \textbf{before}. That reversal is the only thing to drill, and it follows from the rule that has held since the first sentence you read: \textbf{the verb goes last, and everything else queues in front of it.}
\end{grammarlens}

\subsection*{Your turn}
Say these aloud:

\begin{itemize}
\raggedright
  \item \textquotedblleft{}\hi{आदमी} \hi{बच्चे} \hi{से} \hi{बड़ा} \hi{है}\textquotedblright{}
  \item \textquotedblleft{}\hi{बच्चा} \hi{आदमी} \hi{से} \hi{छोटा} \hi{है}\textquotedblright{} — the same fact from the other end
  \item what happens to the adjective in a Hindi comparison (nothing)
  \item which side of the adjective the compared thing sits on
\end{itemize}

\begin{tcolorbox}[breakable,colback=teal!4,colframe=teal!35!black,title={Before you move on}]
What marks the thing compared against? (\textbf{\hi{से}}.) Does the adjective change? (No.) Where does the compared thing stand, against English? (Before the adjective rather than after it.)
\end{tcolorbox}
\section[kitnā]{\hi{कितना} — a word let out of one sentence}

\label{lesson:HI-C92-kitna}



\begin{quote}
Say the question you learned for asking somebody's age. One word in it means something on its own, and you have never used it anywhere else.
\end{quote}

\subsection*{You'll want to know: \hi{कितना}}
\textbf{\hi{कितना}} (\emph{kitnā}) — how much, how many.

\textbf{\hi{कितनी} \hi{किताबें}?} — How many books?

\begin{cousinweb}
You met it inside the age question, where it was one syllable of a memorised lump. It is a word, and it behaves like every \textbf{-\hi{आ}} word in this language:

\noindent
\begin{tabularx}{\linewidth}{@{}>{\raggedright\arraybackslash}X>{\raggedright\arraybackslash}X@{}}
\toprule
\textbf{what is being counted} & \textbf{the form} \\
\midrule
a masculine thing & \hi{कितना} \\
a feminine thing, or a plural & \hi{कितनी} \\
\bottomrule
\end{tabularx}

\textbf{\hi{कितनी} \hi{किताबें}} is feminine and plural, so the question word bends to match its noun — which is the same agreement you have now seen on adjectives, verbs, the past copula, the continuous and the hedge.

\textbf{Hindi does not split \emph{how much} from \emph{how many}.} English needs two phrases because it sorts nouns into countable and uncountable; Hindi asks one question and lets the noun answer.
\end{cousinweb}

\begin{grammarlens}[title={Grammar Lens: the shape of a quantity question}]
\noindent
\begin{tabularx}{\linewidth}{@{}>{\raggedright\arraybackslash}X>{\raggedright\arraybackslash}X@{}}
\toprule
\textbf{question} & \textbf{answer} \\
\midrule
\hi{कितनी} \hi{किताबें}? & \hi{दो} \hi{किताबें।} \\
\hi{कितनी} \hi{हिंदी}? & \hi{थोड़ी।} \\
\bottomrule
\end{tabularx}

The second pair is the one a listening paper sets: a quantity question, and a hedge for an answer. Both words in that exchange bend to the same feminine noun, which is why they rhyme.
\end{grammarlens}

\subsection*{Your turn}
Say these aloud:

\begin{itemize}
\raggedright
  \item \textquotedblleft{}\hi{कितनी} \hi{किताबें}?\textquotedblright{} — how many books?
  \item an answer with a number
  \item \textquotedblleft{}\hi{कितनी} \hi{हिंदी}?\textquotedblright{} and answer it with the hedge
  \item the age question, and name the word inside it you now own
\end{itemize}

\begin{tcolorbox}[breakable,colback=teal!4,colframe=teal!35!black,title={Before you move on}]
How many words does Hindi use for \emph{how much} and \emph{how many}? (One.) Why is it \textbf{\hi{कितनी}} in front of \textbf{\hi{किताबें}}? (The noun is feminine and plural.) Where had you already been saying it? (Inside the age question.)
\end{tcolorbox}
\section[kitnā acchā!]{\hi{कितना} \hi{अच्छा}! — a question that has stopped asking}

\label{lesson:HI-C92-kitna-achchha}



\begin{quote}
You have \textbf{\hi{कितना}} for \emph{how much} and \textbf{\hi{अच्छा}} for \emph{good}. Put them together and listen to what the pair wants to be.
\end{quote}

\subsection*{You'll want to know: \hi{कितना} \hi{अच्छा}!}
\textbf{\hi{कितना} \hi{अच्छा}!} — How good! How lovely!

\textbf{\hi{कितनी} \hi{अच्छी} \hi{किताब}!} — What a good book!

\begin{cousinweb}
Nothing has been added. This is the question word in front of an adjective, with the question taken out of it.

\noindent
\begin{tabularx}{\linewidth}{@{}>{\raggedright\arraybackslash}X>{\raggedright\arraybackslash}X@{}}
\toprule
\textbf{what you say} & \textbf{what it is} \\
\midrule
\hi{कितना} \hi{अच्छा} \hi{है}? & how good is it? \\
\hi{कितना} \hi{अच्छा}! & how good! \\
\bottomrule
\end{tabularx}

The difference is the verb and the voice. Drop \textbf{\hi{है}}, let the pitch fall at the end, and a question becomes an exclamation. English does the identical thing — \emph{how good is it?} against \emph{how good!} — so there is nothing new to translate, only a habit to notice.
\end{cousinweb}

\begin{grammarlens}[title={Grammar Lens: it still agrees, both words at once}]
\noindent
\begin{tabularx}{\linewidth}{@{}>{\raggedright\arraybackslash}X>{\raggedright\arraybackslash}X@{}}
\toprule
\textbf{what is being admired} & \textbf{what you say} \\
\midrule
a masculine thing & \hi{कितना} \hi{अच्छा}! \\
a feminine thing & \hi{कितनी} \hi{अच्छी}! \\
\bottomrule
\end{tabularx}

\textbf{Both} words move together, because both are \textbf{-\hi{आ}} words and both are agreeing with the same noun. That is why the phrase rhymes with itself in either gender, and the rhyme is a useful check: if the two endings do not match, one of them is wrong.

This is the commonest way a Hindi speaker pays a compliment, and it costs you nothing you did not already have.
\end{grammarlens}

\subsection*{Your turn}
Say these aloud:

\begin{itemize}
\raggedright
  \item \textquotedblleft{}\hi{कितना} \hi{अच्छा}!\textquotedblright{} — how good!
  \item \textquotedblleft{}\hi{कितनी} \hi{अच्छी} \hi{किताब}!\textquotedblright{} — what a good book!
  \item the question form of the same phrase
  \item what you take away to turn the question into the exclamation
\end{itemize}

\begin{tcolorbox}[breakable,colback=teal!4,colframe=teal!35!black,title={Before you move on}]
What is removed to make the exclamation? (The verb, and the rising voice.) How many words bend in \textbf{\hi{कितनी} \hi{अच्छी}}? (Both.) What does the rhyme tell you? (That the two endings agree — if they do not match, one is wrong.)
\end{tcolorbox}
\section[mātrā kī tālikā]{mātrā kī tālikā — degree, on one page}

\label{lesson:HI-C92-repaso-maatra}



\begin{quote}
Name the word that hedges an answer and the word that asks a quantity, and say what the two have in common.
\end{quote}

\begin{grammarlens}[title={Grammar Lens: the five}]
\noindent
\begin{tabularx}{\linewidth}{@{}>{\raggedright\arraybackslash}X>{\raggedright\arraybackslash}X@{}}
\toprule
\textbf{what you want to do} & \textbf{what you say} \\
\midrule
hedge an answer & \hi{थोड़ा} / \hi{थोड़ी} \\
describe how an action went & \hi{अच्छा} / \hi{बुरा} \\
measure one thing against another & X \hi{से} \hi{बड़ा} \\
ask a quantity & \hi{कितना} / \hi{कितनी} \\
admire something & \hi{कितना} \hi{अच्छा}! \\
\bottomrule
\end{tabularx}

\textbf{Every one of the five ends in -\hi{आ} and bends.} That is the thread through the chapter, and it means none of them is a new kind of word — they are adjectives doing five different jobs.
\end{grammarlens}

\begin{grammarlens}[title={Grammar Lens: what Hindi does without}]
\noindent
\begin{tabularx}{\linewidth}{@{}>{\raggedright\arraybackslash}X>{\raggedright\arraybackslash}X@{}}
\toprule
\textbf{English has} & \textbf{Hindi has} \\
\midrule
good \textbf{and} well & one word \\
how much \textbf{and} how many & one word \\
bigger, more interesting & the adjective, unchanged \\
\bottomrule
\end{tabularx}

Three places where English carries extra machinery and Hindi carries none. A learner coming the other way has to build those distinctions; you get to drop them.
\end{grammarlens}

\begin{grammarlens}[title={Grammar Lens: \hi{से}, doing a third job}]
\noindent
\begin{tabularx}{\linewidth}{@{}>{\raggedright\arraybackslash}X>{\raggedright\arraybackslash}X@{}}
\toprule
\textbf{use} & \textbf{example} \\
\midrule
from & \hi{घर} \hi{से} \\
by way of & \hi{बस} \hi{से} \\
measured against & \hi{आदमी} \hi{से} \hi{बड़ा} \\
\bottomrule
\end{tabularx}

One postposition, three jobs, and the same idea under all three: the thing \textbf{proceeds from} the noun. A comparison starts at the thing you measure against.
\end{grammarlens}

\subsection*{Your turn}
Say these aloud:

\begin{itemize}
\raggedright
  \item all five rows of the first table
  \item \textquotedblleft{}\hi{आदमी} \hi{बच्चे} \hi{से} \hi{बड़ा} \hi{है}\textquotedblright{}
  \item \textquotedblleft{}\hi{कितनी} \hi{अच्छी} \hi{किताब}!\textquotedblright{}
  \item the three jobs \textbf{\hi{से}} now does
\end{itemize}

\begin{tcolorbox}[breakable,colback=teal!4,colframe=teal!35!black,title={Before you move on}]
What do all five words in this chapter have in common? (They end in \textbf{-\hi{आ}} and bend.) Does a Hindi adjective change to make a comparison? (No.) How many jobs does \textbf{\hi{से}} now do? (Three.)
\end{tcolorbox}
\section[hā̃, thoṛī]{\hi{हाँ}, \hi{थोड़ी} — two words, and one of them is the test}

\label{lesson:HI-C92-sintesis-maatra}



\begin{quote}
You have \textbf{\hi{हाँ}} for yes and the hedge from this chapter. Put them together and you have a complete answer.
\end{quote}

\subsection*{The exchange}
\begin{quote}
— \hi{क्या} \hi{आप} \hi{हिंदी} \hi{बोलती} \hi{हैं}? — \hi{हाँ}, \hi{थोड़ी।}
\end{quote}

That answer is a listening item on the first practice paper. Two words, and \textbf{the second one carries the whole thing}: a bare \emph{hā̃} would claim fluency, and \emph{thoṛī} is what makes the answer true.

The ending is not decoration either. It is \textbf{\hi{थोड़ी}} and not \textbf{\hi{थोड़ा}} because \emph{hindī} is feminine — so a listener who hears the wrong ending hears a speaker who has not noticed what they are talking about.

\subsection*{How to answer}
Say each of these out loud:

\begin{itemize}
\raggedright
  \item yes, a little
  \item I speak a little Hindi
  \item he speaks well
  \item what a good book!
\end{itemize}

Every one of the four ends in a word that had to agree with something. Say them again and listen for the ending each time rather than for the word.

\begin{grammarlens}[title={Grammar Lens: the chapter in one sentence}]
\noindent
\begin{tabularx}{\linewidth}{@{}>{\raggedright\arraybackslash}X>{\raggedright\arraybackslash}X@{}}
\toprule
\textbf{what a paper asks} & \textbf{what it is testing} \\
\midrule
how much do you speak? & the hedge, and its ending \\
how well does he speak? & an adjective doing an adverb's job \\
which is bigger? & \textbf{\hi{से}}, and an unchanged adjective \\
\bottomrule
\end{tabularx}

Three questions, and none of them needs a word you did not already half have. \textbf{This chapter added five words and no new machinery.}
\end{grammarlens}

\subsection*{Your turn}
Say these aloud:

\begin{itemize}
\raggedright
  \item \textquotedblleft{}\hi{हाँ}, \hi{थोड़ी}\textquotedblright{}
  \item why the ending is \textbf{-\hi{ई}} and not \textbf{-\hi{आ}}
  \item \textquotedblleft{}\hi{वह} \hi{अच्छा} \hi{बोलता} \hi{है}\textquotedblright{}
  \item \textquotedblleft{}\hi{आदमी} \hi{बच्चे} \hi{से} \hi{बड़ा} \hi{है}\textquotedblright{}
\end{itemize}

\begin{tcolorbox}[breakable,colback=teal!4,colframe=teal!35!black,title={Before you move on}]
Which of the two words in the answer carries the meaning? (\textbf{\hi{थोड़ी}} — \emph{hā̃} alone would overclaim.) Why that ending? (\emph{Hindī} is feminine.) What marks the thing a comparison measures against? (\textbf{\hi{से}}.)
\end{tcolorbox}
